\documentclass[12pt,a4paper]{article}

\usepackage{fontawesome5}
\providecommand{\faPencilAlt}{\faIcon{pencil-alt}}

\usepackage{styles/couleurs}
\usepackage{styles/mise_en_page}
\usepackage{styles/boites}
\usepackage{styles/maths}
\usepackage{styles/figures}

\usepackage{enumitem}
\usepackage{array}
\usepackage{tabularx}
\usepackage{amsmath}
\usepackage{amssymb}
\usepackage{graphicx}
\usepackage{tikz}
\usepackage{xparse}

\usetikzlibrary{calc,arrows.meta}

\setlist[itemize]{leftmargin=1.4em,itemsep=0.15em,topsep=0.2em}
\setlist[enumerate]{leftmargin=*,itemsep=0.22em,topsep=0.25em}

\newcommand{\lignepointillee}{%
  \noindent\dotfill\par
  \vspace{0.45em}
}

\newcommand{\lignesreponse}[1]{%
  \begin{tcolorbox}[
    enhanced,
    colback=white,
    colframe=NeutreC!40,
    boxrule=0.5pt,
    arc=1.5mm,
    left=3mm,
    right=3mm,
    top=2mm,
    bottom=1mm,
    before skip=2pt,
    after skip=2pt
  ]
    \foreach \i in {1,...,#1}{\lignepointillee}
  \end{tcolorbox}
}

\newcommand{\imageactivite}[3]{%
  \IfFileExists{#3}{%
    \includegraphics[width=#1,height=#2,keepaspectratio]{#3}%
  }{%
    \fbox{%
      \parbox[c][#2][c]{#1}{%
        \centering\sffamily\small
        Image à intégrer\par
        \vspace{0.3em}
        \texttt{\detokenize{#3}}
      }%
    }%
  }%
}

\newcommand{\titreseq}[1]{%
\begin{center}
  {\sffamily\LARGE\bfseries\color{TitreC}
  Séquence 5 -- Statistiques et proportionnalité}\\[0.2em]
  {\sffamily\large\bfseries\color{SavoirC}#1}
\end{center}
}

\begin{document}

\titreseq{Activité 5 -- D'une enquête aux fréquences}

On interroge \(40\) élèves sur leur mode de transport principal pour venir au collège.

\[
\begin{array}{c|cccc}
\text{Mode de transport} & \text{Bus} & \text{À pied} & \text{Voiture} & \text{Vélo}\\
\hline
\text{Effectif} & 16 & 10 & 8 & 6
\end{array}
\]

\section*{Partie A -- Vocabulaire}

Complète.

\begin{tcolorbox}[colback=SavoirC!5,colframe=SavoirC!60,boxrule=0.7pt,arc=2mm]
La population étudiée est : \dotfill

Le caractère étudié est : \dotfill

Les modalités sont : \dotfill

L'effectif total est : \dotfill
\end{tcolorbox}

\section*{Partie B -- Fréquences et pourcentages}

Complète le tableau.

\renewcommand{\arraystretch}{1.5}
\begin{tabularx}{\textwidth}{
|>{\raggedright\arraybackslash}p{3cm}
|>{\centering\arraybackslash}X
|>{\centering\arraybackslash}X
|>{\centering\arraybackslash}X
|>{\centering\arraybackslash}X
|>{\centering\arraybackslash}X|}
\hline
 & \textbf{Bus} & \textbf{À pied} & \textbf{Voiture} & \textbf{Vélo} & \textbf{Total}\\
\hline
\textbf{Effectif} & 16 & 10 & 8 & 6 & \\
\hline
\textbf{Fréquence} & & & & & \\
\hline
\textbf{Pourcentage} & & & & & \\
\hline
\end{tabularx}

\section*{Partie C -- Interpréter}

Réponds par une phrase.

\begin{enumerate}
  \item Quel est le mode de transport le plus fréquent ?
  \item Quelle proportion des élèves vient à pied ?
  \item Combien d'élèves viennent autrement qu'en bus ?
  \item Vérifie que la somme des fréquences vaut \(1\) et que celle des pourcentages vaut \(100\%\).
\end{enumerate}

\lignesreponse{6}

\section*{Partie D -- Comparer}

Dans une autre classe, \(30\%\) des élèves viennent en bus.
Peut-on affirmer qu'il y a moins d'élèves qui viennent en bus que dans cette classe ?
Explique pourquoi il faut connaître l'effectif total.

\lignesreponse{4}

\end{document}
